\documentclass[11pt,letterpaper]{article}
% Reusable style. Build each .tex from its own directory using scripts/build.sh.
\usepackage[T1]{fontenc}
\usepackage[utf8]{inputenc}
\usepackage{lmodern}
\usepackage[margin=0.86in,headheight=15pt]{geometry}
\usepackage{amsmath,amssymb,amsthm,mathtools}
\usepackage{booktabs,longtable,tabularx,array,multicol}
\usepackage[table]{xcolor}
\usepackage{microtype}
\usepackage{enumitem}
\usepackage{fancyhdr}
\usepackage{titlesec}
\usepackage{hyperref}
\usepackage{listings}
\usepackage[most]{tcolorbox}
\usepackage{needspace}
\definecolor{aeNavy}{HTML}{172D43}
\definecolor{aeBlue}{HTML}{176B90}
\definecolor{aeTeal}{HTML}{148477}
\definecolor{aePale}{HTML}{F2F7F9}
\definecolor{aeGray}{HTML}{5C6771}
\hypersetup{colorlinks=true,linkcolor=aeBlue,urlcolor=aeBlue,pdfauthor={Algorithm Engineering Documentation Suite},pdfsubject={Independent educational materials based on the supplied TAOCP volume topics}}
\titleformat{\section}{\Large\bfseries\color{aeNavy}}{\thesection}{0.6em}{}
\titleformat{\subsection}{\large\bfseries\color{aeBlue}}{\thesubsection}{0.6em}{}
\titleformat{\subsubsection}{\normalsize\bfseries\color{aeNavy}}{\thesubsubsection}{0.6em}{}
\setlist[itemize]{leftmargin=*,topsep=3pt,itemsep=1pt}
\setlist[enumerate]{leftmargin=*,topsep=3pt,itemsep=3pt}
\setlength{\parindent}{0pt}\setlength{\parskip}{5pt}
\pagestyle{fancy}\fancyhf{}\fancyhead[L]{\footnotesize\color{aeGray}ALGORITHM ENGINEERING}\fancyhead[R]{\footnotesize\color{aeGray}DOCUMENTATION SUITE}\fancyfoot[L]{\footnotesize\color{aeGray}Independent study and engineering practice}\fancyfoot[R]{\footnotesize\thepage}
\renewcommand{\headrulewidth}{0.35pt}
\newtcolorbox{engineeringnote}[1][]{breakable,colback=aePale,colframe=aeBlue,boxrule=0.5pt,arc=1mm,fonttitle=\bfseries,title={Engineering note},#1}
\newtcolorbox{evidencebox}[1][]{breakable,colback=white,colframe=aeTeal,boxrule=0.6pt,arc=1mm,fonttitle=\bfseries,title={Evidence and deliverables},#1}
\lstdefinestyle{algo}{basicstyle=\ttfamily\footnotesize,breaklines=true,frame=single,rulecolor=\color{aeGray!40},backgroundcolor=\color{aePale},showstringspaces=false,tabsize=2,columns=fullflexible,keepspaces=true,keywordstyle=\color{aeBlue}\bfseries,commentstyle=\color{aeTeal}}
\lstset{style=algo}
\newcommand{\DocTitle}[3]{%
\hypersetup{pageanchor=false}\begin{titlepage}\vspace*{1.2cm}
{\color{aeTeal}\rule{\linewidth}{2.3pt}}\par\vspace{1.15cm}
{\fontsize{24}{30}\selectfont\bfseries\color{aeNavy}\raggedright #1\par}\vspace{0.45cm}
{\Large\color{aeBlue}#2\par}\vspace{1.3cm}
{\large\bfseries #3\par}\vspace{0.35cm}
{\color{aeGray}Edition 1.0 \quad | \quad 29 October 2026\par}\vfill
{\color{aeGray}\small Original educational material aligned to the chapter and section topics in the five supplied volumes of Donald E. Knuth's \emph{The Art of Computer Programming}. This document is not an authorized reproduction or substitute for the original books.}\par\vspace{0.4cm}
{\color{aeTeal}\rule{\linewidth}{0.75pt}}
\end{titlepage}\hypersetup{pageanchor=true}\pagenumbering{roman}}
\newcommand{\SourceNote}[1]{\par\smallskip{\footnotesize\color{aeGray}\textbf{Source locator:} #1}\par\smallskip}
\newcommand{\Outcome}[1]{\par\textbf{Competency target.} #1\par}
\newcommand{\Checkpoint}[1]{\par\textbf{Checkpoint.} #1\par}
% Intentionally simple printable writing lines: call \AnswerLines{number}.
\newcount\alcount
\newcommand{\AnswerLines}[1]{\par\alcount=0\loop\ifnum\alcount<#1\par\noindent\rule{\linewidth}{0.3pt}\vspace{0.33cm}\advance\alcount by1\repeat\par}
\newcommand{\FormField}[2][0.62\linewidth]{\noindent\textbf{#2}\hspace{0.65em}\rule{#1}{0.4pt}\par\vspace{0.18cm}}
\newcommand{\ReviewTable}{\begin{tabularx}{\linewidth}{@{}p{0.23\linewidth}X p{0.2\linewidth}@{}}\toprule Criterion & Recorded evidence & Rating (0--4)\\\midrule Correctness & & \\[1em] Proof quality & & \\[1em] Complexity & & \\[1em] Testing & & \\[1em] Reproducibility & & \\[1em] Communication & & \\[1em]\bottomrule\end{tabularx}}

\begin{document}
\DocTitle{Volume 4A Study Guide}{Chapter 7, Sections 7.1--7.2.1: Boolean methods and enumeration}{Chapter- and section-aligned study guide}
\tableofcontents
\clearpage
\pagenumbering{arabic}
\section{Study protocol}

For each subsection, read the specified part of the original volume, keep a mathematical notebook, produce an original implementation or design sketch, and evaluate the acceptance gate. The investigations below are new exercises, not copies of the book's problems or solutions. Annotate any starred or mathematically advanced source sections as optional on the first pass; revisit them after the prerequisite mathematics is in place.
\begin{evidencebox}Recommended artifacts: study notes, problem specification, derivation, runnable code or model, tests, performance evidence and reflective comparison to original assumptions.\end{evidencebox}

\section{\S 7.1.1: Boolean basics}
\SourceNote{Knuth, Volume 4A, \S 7.1.1.}
\Outcome{Use truth tables, normal forms and symbolic manipulation for compact finite-state computation.}
\subsection{Independent investigation}
Build truth tables for equivalence and implication; validate transformations exhaustively.
\subsection{Analysis and verification gate}
Distinguish logical equivalence from equal runtime cost.
\subsection{Study questions}
\begin{enumerate}
\item Which representations, input assumptions and cost metrics control this technique's applicability?
\item Which invariant or combinatorial identity establishes correctness, and which input is its strongest counterexample candidate?
\item What result would cause the implementation to fail review despite passing a small collection of examples?
\end{enumerate}
\Checkpoint{Preserve a short explanation, the relevant test evidence, and one unresolved question before proceeding.}

\section{\S 7.1.2: Boolean evaluation}
\SourceNote{Knuth, Volume 4A, \S 7.1.2.}
\Outcome{Examine short-circuit evaluation, decision structure and representation.}
\subsection{Independent investigation}
Count evaluation steps for expression trees over different input distributions.
\subsection{Analysis and verification gate}
Show semantics preservation when reordering only under valid assumptions.
\subsection{Study questions}
\begin{enumerate}
\item Which representations, input assumptions and cost metrics control this technique's applicability?
\item Which invariant or combinatorial identity establishes correctness, and which input is its strongest counterexample candidate?
\item What result would cause the implementation to fail review despite passing a small collection of examples?
\end{enumerate}
\Checkpoint{Preserve a short explanation, the relevant test evidence, and one unresolved question before proceeding.}

\section{\S 7.1.3: Bitwise tricks and techniques}
\SourceNote{Knuth, Volume 4A, \S 7.1.3.}
\Outcome{Analyze masks, shifts, bit scans and population count under finite word width.}
\subsection{Independent investigation}
Implement fixed-width bitset operations, validating against set-based operations.
\subsection{Analysis and verification gate}
Test all bits including the top bit and specify signedness behavior.
\subsection{Study questions}
\begin{enumerate}
\item Which representations, input assumptions and cost metrics control this technique's applicability?
\item Which invariant or combinatorial identity establishes correctness, and which input is its strongest counterexample candidate?
\item What result would cause the implementation to fail review despite passing a small collection of examples?
\end{enumerate}
\Checkpoint{Preserve a short explanation, the relevant test evidence, and one unresolved question before proceeding.}

\section{\S 7.1.4: Binary decision diagrams}
\SourceNote{Knuth, Volume 4A, \S 7.1.4.}
\Outcome{Understand reduced ordered BDD representation, sharing and variable-order sensitivity.}
\subsection{Independent investigation}
Construct a small BDD and test two different variable orders against the same truth table.
\subsection{Analysis and verification gate}
Report node counts and prove equivalence of the represented Boolean function.
\subsection{Study questions}
\begin{enumerate}
\item Which representations, input assumptions and cost metrics control this technique's applicability?
\item Which invariant or combinatorial identity establishes correctness, and which input is its strongest counterexample candidate?
\item What result would cause the implementation to fail review despite passing a small collection of examples?
\end{enumerate}
\Checkpoint{Preserve a short explanation, the relevant test evidence, and one unresolved question before proceeding.}

\section{\S 7.2.1.1: Generating all n-tuples}
\SourceNote{Knuth, Volume 4A, \S 7.2.1.1.}
\Outcome{Distinguish lexicographic counting, mixed radix and Gray-like orderings.}
\subsection{Independent investigation}
Write a lazy mixed-radix iterator with explicit bounds.
\subsection{Analysis and verification gate}
Verify cardinality $\prod_i b_i$, no duplicates and termination.
\subsection{Study questions}
\begin{enumerate}
\item Which representations, input assumptions and cost metrics control this technique's applicability?
\item Which invariant or combinatorial identity establishes correctness, and which input is its strongest counterexample candidate?
\item What result would cause the implementation to fail review despite passing a small collection of examples?
\end{enumerate}
\Checkpoint{Preserve a short explanation, the relevant test evidence, and one unresolved question before proceeding.}

\section{\S 7.2.1.2: Generating all permutations}
\SourceNote{Knuth, Volume 4A, \S 7.2.1.2.}
\Outcome{Compare lexicographic next-permutation and swap-based recursive enumeration.}
\subsection{Independent investigation}
Enumerate all permutations for distinct symbols and count generation steps.
\subsection{Analysis and verification gate}
Validate count $n!$, uniqueness and deterministic output ordering.
\subsection{Study questions}
\begin{enumerate}
\item Which representations, input assumptions and cost metrics control this technique's applicability?
\item Which invariant or combinatorial identity establishes correctness, and which input is its strongest counterexample candidate?
\item What result would cause the implementation to fail review despite passing a small collection of examples?
\end{enumerate}
\Checkpoint{Preserve a short explanation, the relevant test evidence, and one unresolved question before proceeding.}

\section{\S 7.2.1.3: Generating all combinations}
\SourceNote{Knuth, Volume 4A, \S 7.2.1.3.}
\Outcome{Implement combinatorial successor and rank-sensitive representation.}
\subsection{Independent investigation}
Enumerate k-subsets of n in sorted order, preserving O(k) output space.
\subsection{Analysis and verification gate}
Compare count to $\binom{n}{k}$ and handle $k=0$ and $k=n$.
\subsection{Study questions}
\begin{enumerate}
\item Which representations, input assumptions and cost metrics control this technique's applicability?
\item Which invariant or combinatorial identity establishes correctness, and which input is its strongest counterexample candidate?
\item What result would cause the implementation to fail review despite passing a small collection of examples?
\end{enumerate}
\Checkpoint{Preserve a short explanation, the relevant test evidence, and one unresolved question before proceeding.}

\section{\S 7.2.1.4--7.2.1.6: Integer/set partitions and trees}
\SourceNote{Knuth, Volume 4A, \S 7.2.1.4--7.2.1.6.}
\Outcome{Understand structure counting and canonical generation to avoid duplicates.}
\subsection{Independent investigation}
Create integer partitions for n up to 8 and verify against an independent count.
\subsection{Analysis and verification gate}
Specify canonical ordering and explain why permutation of summands is not a new partition.
\subsection{Study questions}
\begin{enumerate}
\item Which representations, input assumptions and cost metrics control this technique's applicability?
\item Which invariant or combinatorial identity establishes correctness, and which input is its strongest counterexample candidate?
\item What result would cause the implementation to fail review despite passing a small collection of examples?
\end{enumerate}
\Checkpoint{Preserve a short explanation, the relevant test evidence, and one unresolved question before proceeding.}

\section{Volume synthesis and mastery review}
\begin{enumerate}
\item Draw a dependency graph connecting at least four major ideas from this volume.
\item Select two alternative implementations for one workload and compare their preconditions, proven cost, memory use, measured performance and failure cases.
\item Present a five-minute technical defense of the most surprising experimental result. A claim that differs from the mathematical model requires a measured or documented explanation.
\end{enumerate}
\subsection{Completion checklist}
\begin{itemize}\item All module notes include a source locator and an independently stated learning outcome.
\item Every implementation has reproducible tests and edge cases.
\item Every complexity claim identifies its cost model and underlying assumptions.
\item A peer can independently reproduce one chosen experiment using its recorded configuration.
\end{itemize}
\end{document}
